\documentclass[11pt,a4paper]{article}
\usepackage[utf8]{inputenc}
\usepackage[T1]{fontenc}
\usepackage{lmodern}
\usepackage[a4paper,top=2cm,bottom=2cm,left=2.4cm,right=2.4cm]{geometry}
\usepackage{microtype}
\usepackage{enumitem}
\usepackage{booktabs}
\usepackage[hidelinks]{hyperref}
\usepackage{parskip}
\pagestyle{empty}

\newcommand{\mstitle}{Auditing spatial information leakage in ecological prediction:
quantifying the optimism in map accuracy with the \texttt{spLeakage} R~package}

\begin{document}

\begin{center}
{\large\bfseries Submission metadata}\\[2pt]
\emph{Methods in Ecology and Evolution}\\[6pt]
\textbf{Manuscript:} \mstitle\\
\textbf{Author:} Olatunji Johnson (University of Manchester)
\end{center}

\hrule
\bigskip

% ===========================================================================
\section*{Impact statement}

\textbf{Short version} (\textasciitilde30 words, for a tight character limit):

\begin{quote}
Ecological maps routinely report accuracy that is too optimistic because nearby
survey points leak information between training and testing. We provide a free tool to
detect, quantify and correct this overconfidence in any analysis or published map.
\end{quote}

\textbf{Extended version} (\textasciitilde80 words, if a longer field is available):

\begin{quote}
Maps of where species occur, how much carbon a forest stores, or where disease risk
is highest are only as trustworthy as their reported accuracy---yet that accuracy is
routinely overstated, because spatially close survey points leak information between
training and test data. This paper introduces an open-source R package that takes any
existing analysis, or even a published map, and reports how inflated its accuracy is,
where the leakage lies, and how to correct it---showing across 16 real datasets that
reported accuracy is optimistic every time.
\end{quote}

% ===========================================================================
\section*{Suggested reviewers}

The following researchers have directly relevant expertise in spatial
cross-validation, map-accuracy assessment and species distribution modelling. They are
chosen to span both perspectives the manuscript reconciles---the
nearest-neighbour-matching / spatial-CV view and the design-based view---so the
reconciliation is judged fairly. \emph{(Please verify current affiliation and email at
submission; none is a co-author and there are no collaborative conflicts of
interest.)}

\begin{enumerate}[leftmargin=1.4em,itemsep=6pt]
\item \textbf{Prof.\ Hanna Meyer} --- University of M\"unster, Germany.\\
Lead developer of \texttt{CAST}, the area-of-applicability method and NNDM/kNNDM;
the closest expert on the methods this paper builds on and audits.

\item \textbf{Prof.\ Alexandre M.\ J.-C.\ Wadoux} --- INRAE / LISAH, Montpellier,
France.\\
Leading proponent of design-based map-accuracy assessment; ideal to test that the
paper's reconciliation treats the design-based (``spatial CV is pessimistic'') case
fairly.

\item \textbf{Prof.\ Alexander Brenning} --- Friedrich Schiller University Jena,
Germany.\\
Author of \texttt{sperrorest} and long-standing work on spatial cross-validation and
spatial prediction error; strong on the statistical machinery.

\item \textbf{Dr.\ Roozbeh Valavi} --- (SDM / spatial modelling; \texttt{blockCV}
author).\\
First author of \texttt{blockCV} and of the presence/absence dataset used in the lead
case study; expert on spatial blocking for species distribution models.

\item \textbf{Prof.\ Damaris Zurell} --- University of Potsdam, Germany.\\
Leads work on species-distribution-model evaluation and reporting standards (ODMAP);
well placed to judge the paper's relevance to ecological reporting practice.

\item \textbf{Prof.\ Gurutzeta Guillera-Arroita} --- (quantitative ecology / SDM
methods).\\
Expert on species distribution model evaluation and detectability; a methods-focused
ecological reviewer independent of the spatial-CV development groups.
\end{enumerate}

\textbf{Note on balance.} If the editor prefers to limit reviewers who are themselves
authors of the specific methods extended here (Meyer for NNDM/AOA; Valavi for
\texttt{blockCV}), Brenning, Wadoux, Zurell and Guillera-Arroita provide a
fully independent yet equally expert panel.

\medskip
\textbf{Opposed reviewers.} None.

\end{document}
